\chapter{文档定位与适用范围}

\section{文档目的}
本规范定义课程实验的公开契约：学生工程必须提供的外部行为、接口、资源边界、错误语义与验收条件。它不规定学生的内核目录结构、私有数据结构或具体算法；在本规范未限制处，学生可自行设计实现。

\section{适用对象与前提}
课程面向已学习操作系统概念和 C 语言、但没有操作系统开发经验的学生。课程工程、固定版本工具链、QEMU 配置与 VMware 虚拟机由课程组提供。

\section{术语与约定}
\begin{description}
  \item[必须] 不满足即不能通过对应验收。
  \item[应当] 正常实现应满足；偏离时须在设计说明中解释。
  \item[不要求] 不属于本版验收范围，不得作为隐含评分条件。
\end{description}

\section{文档之间的关系}
本规范优先于实验指导书；指导书提供学习和实现路线，不新增本规范之外的验收要求。测试仅检查已公开的接口和语义。
